\chapter{Context and Requirements}

\section{Problem and Use Cases}
\textit{[Describe the behavior or engineering problem addressed by this RTL module, its users, and the operating scenarios in scope.]}

\section{Module Boundary and Interfaces}
\textit{[List ports, widths, direction, clock domains, reset polarity and behavior, handshake rules, and dependencies on the surrounding design. State what belongs outside this module.]}

% Optional 2:1 layout: keep the interface timing image on the left.
% \HalfModuleDiagram{interface_timing.pdf}
%   {Explain the handshake, timing assumptions, or interface constraint shown.}
%   {Interface timing and handshake conditions.}
%   {fig:interface-timing}

\section{Requirements and Acceptance Criteria}
Give each requirement a stable ID so design decisions and tests can refer to it. Make acceptance criteria observable and measurable.

\begin{tabularx}{\textwidth}{@{}p{17mm}X X@{}}
  \toprule
  ID & Requirement / constraint & Acceptance criterion \\
  \midrule
  REQ-01 & \textit{[Enter requirement]} & \textit{[Metric, threshold, and test conditions]} \\
  \bottomrule
\end{tabularx}

\section{Assumptions and References}
\textit{[Record protocol assumptions, parameter values, target clock, reset expectations, external IP or interface specifications, and sources actually used. Add cited sources to the bibliography.]}
